% GENERATED FILE. Edit canonical lessons, then run npm run generate:books.
% canonical-source-hash: fnv1a64:8bb1ebb735a140de
% canonical-lessons: TE-C359-samrajyam, TE-C359-cakravarti, TE-C359-rakumarudu, TE-C359-rakumari, TE-C359-sthapincu

\chapter{Empires}
\label{ch:te-empires}
\hlchaptermodality{359}

\begin{chapteropening}
\textbf{By the end of this chapter:} \emph{I can say an empire, an emperor, a prince, a princess and to found.}

Complete the last lesson of chapter 359: I can say an empire, an emperor, a prince, a princess and to found.
\end{chapteropening}

\section[\texorpdfstring{\te{సామ్రాజ్యం} (sāmrājyaṁ)}{సామ్రాజ్యం (sāmrājyaṁ)}]{\te{సామ్రాజ్యం} (sāmrājyaṁ) — an empire}

\label{lesson:TE-C359-samrajyam}



\begin{quote}
Before the new one: say the Telugu for to dream of, then the Telugu for luckily.
\end{quote}

\subsection*{You'll want to know: \te{సామ్రాజ్యం}}
\textbf{\te{సామ్రాజ్యం}} — \emph{sāmrājyaṁ} — \textquotedblleft{}an empire\textquotedblright{}.

\begin{quote}
\te{ఆ} \te{సామ్రాజ్యం} \te{సముద్రం} \te{వరకు} \te{ఉండేది}.
\end{quote}

\begin{quote}
\emph{ā sāmrājyaṁ samudraṁ varaku uṇḍēdi.}
\end{quote}

\begin{quote}
That empire stretched as far as the sea.
\end{quote}

Rule over all: built like \te{రాజ్యం}, a kingdom.

\begin{cousinweb}
\textbf{\te{సామ్రాజ్యం}} (\emph{sāmrājyaṁ}) is borrowed from Sanskrit \emph{sāmrājya}, \textquotedblleft{}complete sovereignty, an empire\textquotedblright{} (Monier-Williams).
\end{cousinweb}

\begin{grammarlens}[title={What you've built}]
One more word for telling what happened.
\end{grammarlens}

\subsection*{Your turn}
\begin{itemize}
\raggedright
  \item \emph{Say it:} \emph{sāmrājyaṁ}
  \item \emph{Say it:} \emph{sāmrājyaṁ}, once more
  \item \emph{Say it:} \emph{ā sāmrājyaṁ samudraṁ varaku uṇḍēdi.} — That empire stretched as far as the sea.
  \item \emph{Recall:} say the Telugu for to dream of, then the Telugu for luckily, then say \emph{sāmrājyaṁ} again
\end{itemize}

\begin{tcolorbox}[breakable,colback=teal!4,colframe=teal!35!black,title={Before you move on}]
What does \te{సామ్రాజ్యం} mean? (\textquotedblleft{}An empire\textquotedblright{}.) Say it once more.
\end{tcolorbox}
\section[\texorpdfstring{\te{చక్రవర్తి} (cakravarti)}{చక్రవర్తి (cakravarti)}]{\te{చక్రవర్తి} (cakravarti) — an emperor}

\label{lesson:TE-C359-cakravarti}



\begin{quote}
Before the new one: say the Telugu for luckily, then the Telugu for an empire.
\end{quote}

\subsection*{You'll want to know: \te{చక్రవర్తి}}
\textbf{\te{చక్రవర్తి}} — \emph{cakravarti} — \textquotedblleft{}an emperor\textquotedblright{}.

\begin{quote}
\te{ఆ} \te{చక్రవర్తి} \te{పేరు} \te{అందరికీ} \te{తెలుసు}.
\end{quote}

\begin{quote}
\emph{ā cakravarti pēru andarikī telusu.}
\end{quote}

\begin{quote}
Everyone knows that emperor's name.
\end{quote}

Literally one who turns the wheel: \te{చక్రం}, a wheel.

\begin{cousinweb}
\textbf{\te{చక్రవర్తి}} (\emph{cakravarti}) is borrowed from Sanskrit \emph{cakravartin}, \textquotedblleft{}a universal ruler\textquotedblright{} (CDIAL 4549), built on \emph{cakra}, \textquotedblleft{}a wheel\textquotedblright{}.
\end{cousinweb}

\begin{grammarlens}[title={What you've built}]
One more word for telling what happened.
\end{grammarlens}

\subsection*{Your turn}
\begin{itemize}
\raggedright
  \item \emph{Say it:} \emph{cakravarti}
  \item \emph{Say it:} \emph{cakravarti}, once more
  \item \emph{Say it:} \emph{ā cakravarti pēru andarikī telusu.} — Everyone knows that emperor's name.
  \item \emph{Recall:} say the Telugu for luckily, then the Telugu for an empire, then say \emph{cakravarti} again
\end{itemize}

\begin{tcolorbox}[breakable,colback=teal!4,colframe=teal!35!black,title={Before you move on}]
What does \te{చక్రవర్తి} mean? (\textquotedblleft{}An emperor\textquotedblright{}.) Say it once more.
\end{tcolorbox}
\section[\texorpdfstring{\te{రాకుమారుడు} (rākumāruḍu)}{రాకుమారుడు (rākumāruḍu)}]{\te{రాకుమారుడు} (rākumāruḍu) — a prince}

\label{lesson:TE-C359-rakumarudu}



\begin{quote}
Before the new one: say the Telugu for an empire, then the Telugu for an emperor.
\end{quote}

\subsection*{You'll want to know: \te{రాకుమారుడు}}
\textbf{\te{రాకుమారుడు}} — \emph{rākumāruḍu} — \textquotedblleft{}a prince\textquotedblright{}.

\begin{quote}
\te{రాకుమారుడు} \te{గుర్రం} \te{మీద} \te{అడవికి} \te{వెళ్ళాడు}.
\end{quote}

\begin{quote}
\emph{rākumāruḍu gurraṁ mīda aḍaviki veḷḷāḍu.}
\end{quote}

\begin{quote}
The prince rode into the forest on a horse.
\end{quote}

Plural: \te{రాకుమారులు}, \emph{rākumārulu}.

\te{రా}-, from \te{రాజ}, a king, and \te{కుమారుడు}, a son.

\begin{cousinweb}
\textbf{\te{రాకుమారుడు}} (\emph{rākumāruḍu}) is built from Sanskrit \emph{rājan}, \textquotedblleft{}a king\textquotedblright{} (CDIAL 10679), and \emph{kumāra}, \textquotedblleft{}a boy; a prince\textquotedblright{} (CDIAL 3303).
\end{cousinweb}

\begin{grammarlens}[title={What you've built}]
One more word for telling what happened.
\end{grammarlens}

\subsection*{Your turn}
\begin{itemize}
\raggedright
  \item \emph{Say it:} \emph{rākumāruḍu}
  \item \emph{Say it:} \emph{rākumāruḍu}, once more
  \item \emph{Say it:} \emph{rākumāruḍu gurraṁ mīda aḍaviki veḷḷāḍu.} — The prince rode into the forest on a horse.
  \item \emph{Recall:} say the Telugu for an empire, then the Telugu for an emperor, then say \emph{rākumāruḍu} again
\end{itemize}

\begin{tcolorbox}[breakable,colback=teal!4,colframe=teal!35!black,title={Before you move on}]
What does \te{రాకుమారుడు} mean? (\textquotedblleft{}A prince\textquotedblright{}.) Say it once more.
\end{tcolorbox}
\section[\texorpdfstring{\te{రాకుమారి} (rākumāri)}{రాకుమారి (rākumāri)}]{\te{రాకుమారి} (rākumāri) — a princess}

\label{lesson:TE-C359-rakumari}



\begin{quote}
Before the new one: say the Telugu for an emperor, then the Telugu for a prince.
\end{quote}

\subsection*{You'll want to know: \te{రాకుమారి}}
\textbf{\te{రాకుమారి}} — \emph{rākumāri} — \textquotedblleft{}a princess\textquotedblright{}.

\begin{quote}
\te{రాకుమారి} \te{తోటలో} \te{పాట} \te{పాడింది}.
\end{quote}

\begin{quote}
\emph{rākumāri tōṭalō pāṭa pāḍindi.}
\end{quote}

\begin{quote}
The princess sang a song in the garden.
\end{quote}

The feminine of \te{రాకుమారుడు}: -\te{ఇ} for -\te{ఉడు}.

\begin{cousinweb}
\textbf{\te{రాకుమారి}} (\emph{rākumāri}) is built from Sanskrit \emph{rājan}, \textquotedblleft{}a king\textquotedblright{} (CDIAL 10679), and \emph{kumārī}, \textquotedblleft{}a girl\textquotedblright{}, the feminine of \emph{kumāra} (CDIAL 3303).
\end{cousinweb}

\begin{grammarlens}[title={What you've built}]
One more word for telling what happened.
\end{grammarlens}

\subsection*{Your turn}
\begin{itemize}
\raggedright
  \item \emph{Say it:} \emph{rākumāri}
  \item \emph{Say it:} \emph{rākumāri}, once more
  \item \emph{Say it:} \emph{rākumāri tōṭalō pāṭa pāḍindi.} — The princess sang a song in the garden.
  \item \emph{Recall:} say the Telugu for an emperor, then the Telugu for a prince, then say \emph{rākumāri} again
\end{itemize}

\begin{tcolorbox}[breakable,colback=teal!4,colframe=teal!35!black,title={Before you move on}]
What does \te{రాకుమారి} mean? (\textquotedblleft{}A princess\textquotedblright{}.) Say it once more.
\end{tcolorbox}
\section[\texorpdfstring{\te{స్థాపించు} (sthāpiñcu)}{స్థాపించు (sthāpiñcu)}]{\te{స్థాపించు} (sthāpiñcu) — to found; to establish}

\label{lesson:TE-C359-sthapincu}



\begin{quote}
Before the new one: say the Telugu for a prince, then the Telugu for a princess.
\end{quote}

\subsection*{You'll want to know: \te{స్థాపించు}}
\textbf{\te{స్థాపించు}} — \emph{sthāpiñcu} — \textquotedblleft{}to found; to establish\textquotedblright{}.

\begin{quote}
\te{ఆ} \te{రాజు} \te{ఒక} \te{కొత్త} \te{పట్టణం} \te{స్థాపించాడు}.
\end{quote}

\begin{quote}
\emph{ā rāju oka kotta paṭṭaṇaṁ sthāpiñcāḍu.}
\end{quote}

\begin{quote}
That king founded a new town.
\end{quote}

Past (I …): \te{స్థాపించాను}, \emph{sthāpiñcānu}.

Also of schools and companies: \te{బడి} \te{స్థాపించు}, to found a school.

\begin{cousinweb}
\textbf{\te{స్థాపించు}} (\emph{sthāpiñcu}) is built on Sanskrit \emph{sthāpayati}, \textquotedblleft{}causes to stand, establishes\textquotedblright{} (CDIAL 13756), with the Telugu verb ending \emph{-iñcu}.
\end{cousinweb}

\begin{grammarlens}[title={What you've built}]
One more word for telling what happened.
\end{grammarlens}

\subsection*{Your turn}
\begin{itemize}
\raggedright
  \item \emph{Say it:} \emph{sthāpiñcu}
  \item \emph{Say it:} \emph{sthāpiñcu}, once more
  \item \emph{Say it:} \emph{ā rāju oka kotta paṭṭaṇaṁ sthāpiñcāḍu.} — That king founded a new town.
  \item \emph{Recall:} say the Telugu for a prince, then the Telugu for a princess, then say \emph{sthāpiñcu} again
\end{itemize}

\begin{tcolorbox}[breakable,colback=teal!4,colframe=teal!35!black,title={Before you move on}]
What does \te{స్థాపించు} mean? (\textquotedblleft{}To found; to establish\textquotedblright{}.) Say it once more.
\end{tcolorbox}
